\documentclass[10pt,a4paper]{amsart}
\usepackage[margin=19mm]{geometry}
\usepackage{amsmath,amssymb,mathtools}
\usepackage[scr=rsfs,cal=euler]{mathalfa}
\usepackage{xfrac}
\usepackage[unicode,hidelinks,bookmarks=false]{hyperref}
\newcommand{\ie}{i.e.\ }
\newcommand{\cf}{cf.\ }
\newcommand{\resp}{respectively\ }
\newcommand{\Subsection}{\subsection}
\providecommand{\nobreakdash}{\nobreak\hbox{-}}
\setlength{\emergencystretch}{2em}
\multlinegap0pt
\usepackage{bm}
\usepackage{amssymb}
\usepackage{enumerate}
\usepackage{float}
\newtheorem{theorem}{Theorem}
\newtheorem{lemma}{Lemma}
\def\E{\mathbb{E}}
\newcommand{\R}{{\mathbb R}}
\title{On the Optimal Control Problem of Stochastic Semilinear Partial Differential Equations with Non-Globally Lipschitz Coefficients}
\author{Oleksiy Kapustyan, Olha Martynyuk, Oleksandr Misiats, Oleksandr Stanzhytskyi}
\date{Source: arXiv:2608.23985v1 (CC BY 4.0)}
\begin{document}
\maketitle

\noindent\emph{Source and licence.} On the Optimal Control Problem of Stochastic Semilinear Partial Differential Equations with Non-Globally Lipschitz Coefficients. Version arXiv:2608.23985v1 (CC BY 4.0), distributed by arXiv under the Creative Commons Attribution 4.0 International licence, http://creativecommons.org/licenses/by/4.0/, as recorded in the arXiv licence metadata for this exact version and shown on the abstract page https://arxiv.org/abs/2608.23985v1. Full licence notice: \texttt{licenses/arxiv-2608.23985-CC-BY-4.0.txt}.\par\smallskip

\noindent\textit{Editorial scope.} Counted unit: the article's theorem Th:2.4 (source0.tex lines 1010--1023). The article's complete original proof of it is delivered with it (source0.tex lines 2914--3139, one \texttt{proof} environment of 226 lines, which the article places away from the statement). No mathematics is written, repaired or re-derived: the statement and the proof are the article's own lines, in the article's own notation, and no retained line is substituted or re-flowed.  Retained uncounted as the source-local context the proof needs: lines 227--237; lines 239--247; lines 249--256; lines 279--287; lines 373--403; lines 990--992; lines 1176--1207; lines 2430--2441.  Nothing was omitted from any retained span, so the package carries no omission record.  No retained line contains a citation, so the wrapper carries no bibliography and no bibliography entry.  No retained line contains a figure environment, an image-inclusion command or drawing code, so no image is delivered and the package declares no asset dependency.  Editorial formatting only, in addition: an AMS \texttt{amsart} preamble that restates the article's own declarations and macros used by the retained lines, namely the environments \texttt{theorem}, \texttt{lemma} on separate counters, together with the standard packages those lines need.  The retained environments print in this document as Theorem 1, Lemma 1 in order of appearance, whereas the article prints the same items with the numbering of its own full document.  The complete native source of the article as distributed in the arXiv e-print for this exact version is bundled unchanged as \texttt{sources/208441/source0.tex}.
\begin{equation}\label{1.1}
\begin{cases}
\begin{aligned}
& dy = \bigl(Ay + f(y) + u(t)\bigr)\,dt
      + \sigma(y)\,dW(t), \\
& y(t,x) = 0,
\qquad x\in \partial D,\; t\in(0,T), \\
& y(0,x) = y_0(x,\omega).
\end{aligned}
\end{cases}
\end{equation}

\begin{equation}\label{1.2}
J(u)
=
\E\int_{0}^{T}\int_{D} y^2(t,x)\,dx\,dt
+
\E\int_{0}^{T}\int_{D} u^2(t,x)\,dx\,dt
\;\to\;
\inf .
\end{equation}

\begin{equation}\label{1.3}
J(u)
=
\E\int_{0}^{\infty}\int_{D} e^{-\gamma t} y^2(t,x)\,dx\,dt
+
\E\int_{0}^{\infty}\int_{D} u^2(t,x)\,dx\,dt
\to \inf 
\end{equation}

\begin{equation}\label{1.4}
\sum_{i,j=1}^{d}
a_{ij}(x)\eta_i\eta_j
\ge
\Lambda_A |\eta|^2,
\qquad
\forall\,\eta\in\mathbb \R^d,
\ \text{for a.e. } x\in D.
\end{equation}

\begin{equation}
\left\{
\begin{aligned}
dy_k
&=
\bigl(Ay_k+f_k(y_k)+u\bigr)\,dt
+\sigma(y_k)\,dW(t),
\\[1ex]
J_k(u)
&=
E\int_0^T\int_D y_k^2(t,x)\,dx\,dt
+
E\int_0^T\int_D u^2(t,x)\,dx\,dt
\to \inf .
\end{aligned}
\right.
\label{1.7}
\end{equation}
 Since the coefficients
$f_k'(s)$ and $\sigma'(s)$ are globally Lipschitz and  possess bounded
derivatives (after a suitable smoothing, if necessary),
 the stochastic maximum principle is applicable to
problem \eqref{1.7}. Loosely speaking, the main result of the present paper is the following: if $
(J^*,u^*,y^*)$ and $
(J_k^*,u_k^*,y_k^*)
$
are the solutions of problems
\eqref{1.1}--\eqref{1.2}
and
\eqref{1.7},
respectively, then:

\begin{equation}\label{2.10}
    \gamma > \mu
\end{equation}

\begin{lemma}\label{lem:3.2}
Suppose the conditions \eqref{1.4} and {{\bf\rm(A1)--(A3)}} hold.
Assume 

\[
 u_n \rightharpoonup u
\qquad\text{weakly in }L^2(\Omega_T;H),
\qquad
n\to\infty,
\]
then, along a subsequence,
\begin{equation} y(t, u_n)
\rightharpoonup
 y(t, u)
\qquad\text{weakly in }
L^2(\Omega_T;H),
\qquad
n\to\infty.
\label{3.1}
\end{equation}
For the ``cut--off'' problem \eqref{1.7} we have a similar result, namely, for every $k \ge 1$,
\begin{equation}
y_k(t, \widetilde u_n)
\rightharpoonup
y_k(t, \widetilde u)
\qquad\text{weakly in }
L^2(\Omega_T;H),
\qquad
n\to\infty.
\label{3.2}
\end{equation}
\end{lemma}

\begin{equation}
\mathbb{E}\|y(t)\|^2
\leq
C_{20}
\left(
\mathbb{E}\|y_0\|_{L^{\infty}}^2 
+
\int_0^t
\mathbb{E}\|u(s)\|^2\,ds
+ \lambda \sup_{i}\|e_i\|_{L^\infty}^2 L_{\sigma}^2 t \right) e^{mt}.
\label{3.35}
\end{equation}

\begin{theorem}\label{Th:2.4}
Assume that conditions \eqref{1.4},
{{\bf \rm(A1)--(A3)}}, and \eqref{2.10} hold.
Then the optimal control problem
\eqref{1.1}--\eqref{1.3}
admits a weak martingale solution

\[
\bigl(
\widetilde u^{\,*}(t),
\widetilde y^{\,*}(t,\widetilde u^{\,*})
\bigr).
\]
\end{theorem}

\begin{proof}[Proof of Theorem \ref{Th:2.4}]
Recall that $U_{\mathrm{ad}}$ denotes the set of all admissible controls
for the problem \eqref{1.1}--\eqref{1.3}. Let
\[
J^*
=
\inf_{u\in U_{\mathrm{ad}}}J(u).
\] 
In view of \eqref{2.10} and
\eqref{3.35},
\[
J(0)
=
\mathbb{E}\int_0^\infty
e^{-\gamma t}
\|y(t,0)\|^2\,dt
<
\infty,
\]
which yields that $0 \in U_{\mathrm{ad}}$. Next, let $\{u_n\}$ be a minimizing sequence for $J$, such that
\begin{equation}
J(u_n)
\to
J^*,
\qquad n\to\infty.
\label{4.18}
\end{equation}
Thus, for sufficiently large $n$,
\begin{equation}
\begin{aligned}
J(u_n)
={}&
\mathbb{E}\int_0^\infty
e^{-\gamma t}
\|y(t,u_n)\|^2\,dt
+
\mathbb{E}\int_0^\infty
\|u_n(t)\|^2\,dt 
\leq
J^*+1.
\end{aligned}
\label{4.19}
\end{equation}
This implies that $\{u_n\}$ is weakly compact in
$L^2(\Omega_T;H)$, and
$\{y(\cdot,u_n)\}$ is weakly compact in
$L_\gamma^2(\Omega_\infty;H)$. Hence, passing to a subsequence if
necessary, we obtain
\begin{equation}
u_n
\rightharpoonup
u
\quad\text{weakly in }
L^2(\Omega_T;H),
\end{equation}
and
\begin{equation}
y(\cdot,u_n)
\rightharpoonup
z(\cdot)
\quad\text{weakly in }
L_\delta^2(\Omega_\infty;H),
\qquad n\to\infty.
\label{4.20}
\end{equation}
Now we show that there exists a probability space
$(\widetilde{\Omega},\widetilde{\mathcal{F}},\widetilde{P})$ such that
\[
\mathcal{L}\bigl(y(\cdot,u_n),u_n\bigr)
=
\mathcal{L}\bigl(\widetilde{y}(\cdot,\widetilde{u}_n),
\widetilde{u}_n\bigr),
\]
\[
\mathcal{L}(z)=\mathcal{L}(\widetilde{z}),
\qquad
\mathcal{L}(u)=\mathcal{L}(\widetilde{u}),
\]
and
\[
\widetilde{z}(t)
=
\widetilde{y}(t,\widetilde{u}).
\]
Indeed, by Lemma \ref{lem:3.2}, for every $N>0$, on $[0,N]$ we have
\begin{equation}
\widetilde{y}(\cdot,\widetilde{u}_n)
\rightharpoonup
\widetilde{y}(\cdot,\widetilde{u})
\quad\text{weakly in }
L^2(\widetilde{\Omega}_N;H),
\qquad n\to\infty.
\end{equation}
Note that the space $(\widetilde{\Omega}, \widetilde{\mathcal{F}}, \widetilde{P})$ may be chosen to be $([0,1], \mathcal{B}, \lambda)$, where $\mathcal{B}$ are Borel subsets of $[0,1]$, and $\lambda$ is Lebesgue measure. Thus, for any
$\varphi\in L_\gamma^2(\Omega_\infty;H)$, we have
\begin{align}
&\left|
\widetilde{\mathbb{E}}
\int_0^\infty
\int_D
\bigl(
\widetilde{y}(t,\widetilde{u}_n)
-
\widetilde{y}(t,\widetilde{u})
\bigr)
\varphi(t,x,\omega)\,dx\,
e^{-\gamma t}\,dt
\right|
\nonumber\\
&\qquad\leq
\left|
\widetilde{\mathbb{E}}
\int_0^N
\int_D
\bigl(
\widetilde{y}(t,\widetilde{u}_n)
-
\widetilde{y}(t,\widetilde{u})
\bigr)
\varphi(t,x,\omega)\,dx\,
e^{-\gamma t}\,dt
\right|
\nonumber\\
&\qquad\quad+
\widetilde{\mathbb{E}}
\int_N^\infty
\int_D
\left|
\widetilde{y}(t,\widetilde{u}_n)
-
\widetilde{y}(t,\widetilde{u})
\right|
\left|
\varphi(t,x,\omega)
\right|\,dx\,
e^{-\gamma t}\,dt
=: I_1+I_2.
\label{4.21}
\end{align}
Since $\varphi\in L_\delta^2(\widetilde{\Omega}_N;H)$,
for every fixed $N$,
\begin{equation}
I_1 \to 0,
\qquad n\to\infty.
\label{4.22}
\end{equation}
As far as the second term in \eqref{4.21}, for some constant $C>0$, we
obtain
\[
I_2
\leq
C
\left(
\int_N^\infty
\left[
\widetilde{\mathbb{E}}
\left\|
\widetilde{y}(t,\widetilde{u}_n)
\right\|^2
+
\widetilde{\mathbb{E}}
\left\|
\widetilde{y}(t,\widetilde{u})
\right\|^2
\right]
e^{-\gamma t}\,dt
\right)^{1/2}
\left(
\int_N^\infty
\widetilde{\mathbb{E}}
\|\varphi(t)\|^2
e^{-\gamma t}\,dt
\right)^{1/2}.
\]
Since $\widetilde{u}_n$ is bounded in
$L^2(\widetilde{\Omega}_\infty;H)$, in conjunction with \eqref{3.35}, we deduce that the first factor in the
right-hand side is uniformly bounded with respect to both $N$ and $n$,
whereas the second factor tends to zero as $N\to\infty$. Thus, for
every $\varepsilon>0$, we may choose $N>0$ sufficiently large so that
\[
I_2<\frac{\varepsilon}{2}
\]
uniformly with respect to $n$. Then, for this fixed $N$, choose
$n$ sufficiently large so that
\[
I_1<\frac{\varepsilon}{2}.
\]
Thus, \eqref{4.20} follows. Finally, using \eqref{4.18} and weak lower semicontinuity of the cost
functional, we obtain
\[
\begin{aligned}
J^*
&=
\lim_{n\to\infty}J(u_n) \geq
\liminf_{n\to\infty}
\widetilde{\mathbb{E}}
\int_0^\infty
\left\|
\widetilde{y}(t,\widetilde{u}_n)
\right\|^2
e^{-\gamma t}\,dt
+
\liminf_{n\to\infty}
\widetilde{\mathbb{E}}
\int_0^\infty
\left\|
\widetilde{u}_n(t)
\right\|^2\,dt
\\
&\geq
\widetilde{\mathbb{E}}
\int_0^\infty
\left\|
\widetilde{y}(t,\widetilde{u})
\right\|^2
e^{-\gamma t}\,dt
+
\widetilde{\mathbb{E}}
\int_0^\infty
\left\|
\widetilde{u}(t)
\right\|^2\,dt.
\end{aligned}
\]
This concludes the proof of Theorem \ref{Th:2.4}.
\end{proof}

\end{document}
